\section{Uniform estimates for selected local observations}
\label{probability:section}

The local observations used later are chosen after inspecting the
environment.  We therefore need estimates for the largest reward over
all admissible choices.  There are two different tail estimates.
Sparse-path counting controls fields with sufficiently light tails,
including bounds obtained by choosing among separated candidates.  For
the tunnel rewards, the available tail is only of order $t^{-2}$;
their control instead uses the thin geometry of the selected diagram
and a directed collection bound.

Selected-path estimates are related to the greedy lattice-animal and
lattice-path bounds studied by Martin \cite{Martin2002}. The critical
quadratic tails here require the additional collection and geometric
deformation arguments proved below.

The collection bound \eqref{probability:collection-assumption} is a
hypothesis, proved for tunnels in Proposition~\ref{access:collection}.
It does not follow from locality and a first-moment estimate.  All
constants below may depend on the fixed lattices, cones, dependence
ranges, and observation templates, but are uniform in the index $n$
of a family of laws.  Constants denoted by $C_F$ may also depend on
the deformation parameter $F$; constants denoted by $C$ do not.

\subsection{Sparse observations and light selected tails}

Use the $\ell^1$ norm on $\Z^2$, and let $\mathcal P_L$ be the
nearest-neighbor paths starting at the origin with at most $L$ edges.
Sums in this subsection count each visited site once.  Allowing at
most $m_0$ occurrences per site multiplies the bounds by $m_0$.

\begin{lemma}[Sparse-path estimate]
\label{probability:sparse}
Let $(I_z)_{z\in\Z^2}$ be indicators with $\PP(I_z=1)\le p$.
Suppose $\Z^2$ has a partition into $q$ classes such that the
variables within each class are mutually independent.  For every
integer $L\ge0$,
\begin{equation}
 \E\max_{\pi\in\mathcal P_L}\sum_{z\in\pi}I_z
 \le C_q\min\{L+1,(L+1)\sqrt p,(L+1)^2p\}.
 \label{probability:sparse-bound}
\end{equation}
The same conclusion, with adjusted constants, holds for lists of
distinct sites whose initial distance from the origin plus the sum
of their successive distances is at most $c(L+1)$, for fixed $c$.
It also holds on any fixed full-rank planar lattice.  In particular,
it applies to a field with a fixed range of dependence.
\end{lemma}

\begin{proof}
First suppose all indicators are independent.  If a path visits
$\ell$ distinct marked sites, list them in their order of first
encounter as $z_1,\ldots,z_\ell$.  With $z_0=0$,
$\sum_{i=1}^{\ell}|z_i-z_{i-1}|_1\le L$.  There are at most
$4(d+1)$ displacements of $\ell^1$ length $d$.  The number of
possible lists is therefore at most
\[
 4^\ell\sum_{d_1+\cdots+d_\ell\le L}
                   \prod_{i=1}^{\ell}(d_i+1)
 =4^\ell\binom{L+2\ell}{2\ell}
 \le \left(\frac{C(L+1)}{\ell}\right)^{2\ell},
 \qquad 1\le\ell\le L+1.
\]
The identity follows from
$\sum_{d\ge0}(d+1)x^d=(1-x)^{-2}$, followed by summing
coefficients through degree $L$.  The count includes lists that
cannot occur along a path, which only increases the bound.

Distinctness makes the probability that a specified list is marked
at most $p^\ell$.  If $M_L$ is the maximum in the statement and
$a=C(L+1)\sqrt p$, the union bound gives
\begin{equation}
 \PP(M_L\ge\ell)
 \le\min\{1,(a/\ell)^{2\ell}\}.
 \label{probability:sparse-union}
\end{equation}
For $a\le1/2$, summing over $\ell\ge1$ gives at most
$\sum_{\ell\ge1}a^{2\ell}\le2a^2$.  For $a>1/2$, split at
$\lceil2a\rceil$: the initial sum is $O(a)$ and the remaining
sum is geometric.  Thus $\E M_L\le C\min\{a,a^2\}$, after
adjusting constants also for $1/2<a<1$.  Combining this with
$M_L\le L+1$ proves \eqref{probability:sparse-bound}.

For a partition into independence classes, apply the same list
count to sites in each class and sum the resulting $q$ maxima.
No independence between classes is needed.  For example, when
$I_z$ depends on independent inputs in $z+[-R,R]^2$, residue
classes modulo an integer greater than $2R$ give the required
partition.  The same partition argument works when sigma-fields
generated by sets at distance greater than a fixed range are
independent.

Only the budget for successive distances entered the counting
argument, so it also proves the assertion about lists.  An invertible
linear identification of a fixed full-rank lattice with $\Z^2$
changes this budget by a fixed factor and proves the lattice version.
\end{proof}

\begin{corollary}[Integrated bound for selected tails]
\label{probability:integrated-tails}
Let $(U_z)$ be nonnegative fields whose dependence range is uniform
in a possible law index $n$, and suppose
$\sup_{n,z}\PP(U_z>t)\le p(t)$.  Then, for $T\ge0$,
\begin{align}
 \E\sup_{\pi\in\mathcal P_L}
       \sum_{z\in\pi}U_z\mathbf1_{\{U_z>T\}}
 &\le C(L+1)\left[T\sqrt{p(T)}+
                         \int_T^\infty\sqrt{p(t)}\,dt\right].
 \label{probability:tail-integral}
\end{align}
In particular, a uniform bound $\PP(U_z>t)\le A t^{-4}$ for
$t\ge1$ gives $C_A(L+1)/T$ for $T\ge1$.  The same bounds hold
for the site lists and bounded multiplicities in
Lemma~\ref{probability:sparse}.
\end{corollary}

\begin{proof}
For $u\ge0$,
\[
 u\mathbf1_{\{u>T\}}
 =T\mathbf1_{\{u>T\}}+
       \int_T^\infty\mathbf1_{\{u>t\}}\,dt.
\]
Take the supremum over paths, move it inside the nonnegative
integral as an upper bound, and apply
Lemma~\ref{probability:sparse} at every level.  Tonelli's theorem
gives \eqref{probability:tail-integral}, possibly with an infinite
right-hand side.  Integrating $\sqrt A\,t^{-2}$ proves the
fourth-power assertion.

The sharper minimum of $(L+1)\sqrt{p(t)}$ and $(L+1)^2p(t)$
in \eqref{probability:sparse-bound} can also be used inside the
integral.  In particular, the estimate has no additive constant
independent of $t$ to be integrated over an infinite interval.
\end{proof}

\begin{lemma}[Selection from separated candidates]
\label{probability:candidates}
A fixed bounded region contains at most $J$ candidate cells.  For
each cell $c$, let $U_c\ge0$ be measurable with respect to a
specified subset of one family of independent inputs, with
$\E U_c\le A$.  Suppose each admissible object provides a tuple
of $m$ cells whose input supports are pairwise disjoint.  There is
a local random variable $V$ that bounds, simultaneously for every
object, its smallest available $U_c$, and satisfies
\[
 \PP(V>t)\le J^m\min\{1,(A/t)^m\},\qquad t>0.
\]
For fixed $J,m,A$, the fourth moment of $V$ is uniformly bounded
when $m>4$.  The same assertion holds conditional on auxiliary
data if the disjoint supports retain their product law and each
conditional first moment is at most the same deterministic $A$.
\end{lemma}

\begin{proof}
Before selecting an object, form
\[
 V=\max_{(c_1,\ldots,c_m)}\min_{1\le i\le m}U_{c_i},
\]
where the maximum ranges over all tuples in the region with
pairwise disjoint supports, including those that no admissible
object supplies.  Set $V=0$ if there are no such tuples.  This is
a local variable and gives the required simultaneous domination.
For each fixed tuple, independence and Markov's inequality bound
the probability that all $m$ values exceed $t$ by
$\min\{1,(A/t)^m\}$.  There are at most $J^m$ tuples, so a
union bound proves the tail estimate.  Finally,
\[
 \E V^4=4\int_0^\infty t^3\PP(V>t)\,dt<\infty
 \qquad(m>4).
\]
The conditional proof is identical.  The object itself may depend
on the whole environment: independence is used only for each
fixed tuple, before taking the finite maximum.
\end{proof}

\subsection{Compressing a thin diagram by a lattice permutation}

The critical-tail argument will compress one coordinate and stretch
the other by reciprocal factors.  The next lemma realizes this
area-preserving deformation as a permutation of the lattice, so
that transporting an iid field preserves its law exactly.

\begin{lemma}[Bounded-displacement matching]
\label{probability:matching}
Let $\Lambda$ be a full-rank lattice in $\R^2$ and let $T$ be
invertible with $|\det T|=1$.  There is a bijection
$\psi:\Lambda\to\Lambda$ and a finite $b_T$ such that
\[
                  |\psi(z)-Tz|\le b_T\qquad(z\in\Lambda).
\]
For a bounded fundamental tile $Q$ of $\Lambda$, one may take
$b_T\le\sup_{q,q'\in Q}|Tq-q'|$.  For an affine map $Tz+a$,
the same bound holds with $|\psi(z)-(Tz+a)|$ on the left.
\end{lemma}

\begin{proof}
Compare the tilings $(T(z+Q))_{z\in\Lambda}$ and
$(w+Q)_{w\in\Lambda}$.  Join $z$ to $w$ when the corresponding
tiles intersect in positive area.  This bipartite graph is countable
and locally finite.  All tiles have the same area $v_Q>0$.  Any union
of $k$ tiles on either side has area $kv_Q$ and is covered, up to a
null set, by its neighboring tiles on the other side.  It therefore
has at least $k$ neighbors.  Both finite Hall conditions hold.

Here is the compactness step giving a perfect matching.  By finite
Hall's theorem~\cite[Theorem~1]{Hall1935}, any prescribed finite
set of left vertices is covered by a finite matching $M_1$, and
any prescribed finite set of right vertices by a finite matching
$M_2$.  The union of the two matchings consists of alternating
paths and cycles, together with common edges.  On each component
one can choose an alternating matching covering all prescribed
vertices.  Internal vertices are covered by either choice.  On an
odd-length path, one choice covers both endpoints.  On an
even-length path, the endpoints lie on the same side, and the
matching that covered the prescribed vertices on that side already
covers every prescribed endpoint.  Thus one matching covers both
prescribed finite sets.

Exhaust the two vertex sets by finite sets and choose such a
matching at each stage.  A diagonal subsequence converges on every
edge.  Since each vertex has finite degree and is covered at all
sufficiently large stages, exactly one of its incident edges belongs
to the limit.  The limit is a perfect matching.

Let $\psi(z)=w$ for the matched pair.  An intersection point has
the form $Tz+Tq=w+q'$, with $q,q'\in Q$, giving the stated
displacement bound.  Translate the first tiling by $a$ to obtain
the affine statement, with the same error relative to $Tz+a$.
\end{proof}

We now specify the geometry and its uniformities.  A cone here is
closed, convex, pointed, and has nonempty interior.  Fix a target
cone $\mathcal K$ strictly inside the positive quadrant, and fix
linearly independent vectors $u_1,u_2$ strictly inside $\mathcal K$,
ordered so that $U=(u_1,u_2)$ has positive determinant.  Let
$E=(e_1,e_2)$ range over a fixed compact set of matrices with
positive determinant.  For each application, $E$ and its lattice
matching are deterministic.  No supremum over an
environment-dependent choice of $E$ is taken.  All constants in
the thin-diagram estimates below are uniform over the allowed
matrices $E$, with $F$-dependence indicated by subscripts.

\begin{lemma}[Thin diagrams and order errors]
\label{probability:thin-geometry}
Fix $H,c>0$ and a full-rank lattice $\Lambda$.  Suppose
$z_1,\ldots,z_k\in\Lambda$ satisfy the all-pairs condition
\begin{equation}
 E^{-1}(z_j-z_i)\ge-H(1,1)
 \quad\hbox{coordinatewise for every }i<j,
 \label{probability:approximate-order}
\end{equation}
and lie in the parallelogram
\begin{equation}
 E\big([-H,cN+H]\times[-H,\delta N+H]\big).
 \label{probability:thin-box}
\end{equation}
For $F\ge1$, define
\[
 T_F=\left(\frac{\det E}{\det U}\right)^{1/2}
          U\begin{pmatrix}F^{-1}&0\\0&F\end{pmatrix}E^{-1}.
\]
Then $\det T_F=1$.  A matching $\psi_F$ supplied by
Lemma~\ref{probability:matching} sends these sites into a
deterministic square of diameter at most
\begin{equation}
             C(N/F+F\delta N)+C_F.
 \label{probability:compressed-size}
\end{equation}
For every fixed $r_0>0$, there is a finite $R_F$, uniform in $E$,
such that, if $i<j$ and $|z_j-z_i|>R_F$, then
$\psi_F(z_j)-\psi_F(z_i)\in\mathcal K$ and both coordinates
of this increment exceed $r_0$.
\end{lemma}

\begin{proof}
The scalar prefactor makes the determinant one.  In $E$-coordinates,
the map contracts the first extent of the box by $F$ and expands
the second by $F$.  The prefactor and $U$ have uniform bounds, so
the transformed diameter is at most
$C(N/F+F\delta N)+C_F$.  The fundamental-tile bound in
Lemma~\ref{probability:matching} bounds the matching error
uniformly in $E$ for fixed $F$.  Adding this error proves
\eqref{probability:compressed-size}.  Its $N$-dependent coefficient
$C$ is independent of $F$; the deformation of the fixed margins
and the matching error are included in $C_F$.

For $v=z_j-z_i$, condition
\eqref{probability:approximate-order} gives $v=Ea+e$ with
$a\ge0$ coordinatewise and $|e|\le CH$.  The vector $T_FEa$
is in the cone generated by $u_1,u_2$, strictly inside
$\mathcal K$.  Consequently its distance to the complement of
$\mathcal K$, and each of its positive coordinates, are at least
$\rho_{\mathcal K}|T_FEa|$ for a fixed $\rho_{\mathcal K}>0$.
For fixed $F$,
\[
 |T_FEa|\ge c_F(|v|-CH),\qquad |T_Fe|\le C_FH,
\]
with $c_F>0$.  Replacing $T_Fv$ by
$\psi_F(z_j)-\psi_F(z_i)$ adds at most twice the matching
error.  Choose $R_F$ so that the cone margin and the two coordinate
margins exceed these errors plus $r_0$ whenever $|v|>R_F$.
Compactness of the allowed matrices makes this choice uniform in
$E$.
\end{proof}

The all-pairs condition is essential: a bounded error at each
successive step may accumulate.  The box condition follows, for
example, when the same all-pairs condition also includes
deterministic initial and terminal points with $E$-coordinates
$(0,0)$ and $(c'N,\delta'N)$, where $0\le c'\le c$ and
$0\le\delta'\le\delta$.  Neither condition bounds the number of
visits to one site; that is a separate hypothesis below.

\subsection{Critical tails in a thin diagram}

Let $(Y_z)_{z\in\Lambda}$ be a nonnegative iid field on a fixed
full-rank lattice.  For $L\ge1$ and $x\in\R^2$, let $G(L,x)$
be the maximum of $\sum_{z\in S}Y_z$ over chains in
$x+[-L,L]^2$ whose successive increments lie in $\mathcal K$
and have both coordinates at least $r_0$.  Include the empty chain
and all singletons.  Suppose the laws, possibly indexed by $n$,
satisfy
\begin{equation}
 \E G_n(L,x)\le C_0L+C_1,
 \qquad L\ge1,\quad x\in\R^2,
 \label{probability:collection-assumption}
\end{equation}
with the same constants for every $n$.  This finite-volume bound
will be proved for the charged tunnels in
Proposition~\ref{access:collection}, before its use in
Section~\ref{links:section}.

\begin{lemma}[Weak second-moment bound]
\label{probability:weak-two}
Under \eqref{probability:collection-assumption},
\[
                \sup_n\PP(Y_z>t)\le C\min\{1,t^{-2}\},
                \qquad t>0.
\]
\end{lemma}

\begin{proof}
Fix a law and $t>0$, and set $p=\PP(Y_z>t)$.  There is nothing
to prove if $p=0$.  Otherwise choose $L$ comparable to $p^{-1/2}$,
with lattice-dependent comparison constants, so that $L\ge1$ and
$[-L,L]^2$ contains at least $p^{-1}$ lattice sites.  Independence
implies that some value in the square exceeds $t$ with probability
at least $1-e^{-1}$.  Since singletons are allowed in $G$,
\[
 (1-e^{-1})t\le\E G_n(L,0)
                  \le C_0L+C_1\le Cp^{-1/2}.
\]
Rearrange and combine with $p\le1$.  The constants do not depend
on the law.
\end{proof}

This is a weak $L^2$ estimate, not an assertion of a finite second
moment.  Substituting its $t^{-2}$ tail into the square-root integral
in \eqref{probability:tail-integral} gives a divergent bound.  The
next argument uses the thin box to control the values that would
cause this divergence.

\begin{lemma}[Negligible selected tails in a thin diagram]
\label{probability:thin-tails}
Assume \eqref{probability:collection-assumption}, and fix
$H,c,m_0$.  For $N\ge1$, let $\mathcal S(N,\delta)$ consist of
finite site sequences satisfying
\eqref{probability:approximate-order} and
\eqref{probability:thin-box}, with at most $m_0$ occurrences per
site.  Suppose also that their distinct sites can be visited by a
path whose initial distance from the origin plus total length is
at most $c(N+1)$.

For every $F\ge1$ there is $C_F<\infty$ such that, for $T\ge1$
and $0\le\delta\le F^{-2}$,
\begin{equation}
 \E\sup_{S\in\mathcal S(N,\delta)}
       \sum_{z\in S}Y_z\mathbf1_{\{Y_z>T\}}
 \le \frac{CN}{F}+\frac{C_F(N+1)}{T}+C_F.
 \label{probability:thin-tail-bound}
\end{equation}
Here $C$ is independent of $F,T,\delta,n$.  Consequently, for
every $\varepsilon>0$ there are $\delta_0>0$ and $T_0<\infty$
such that, for every sequence of laws and every $N(n)\to\infty$,
\[
 \limsup_{n\to\infty}\frac1{N(n)}
 \E\sup_{S\in\mathcal S(N(n),\delta_0)}
       \sum_{z\in S}Y_z\mathbf1_{\{Y_z>T_0\}}
 \le\varepsilon.
\]
Finite sums of field types and independence sublattices satisfy
the same conclusion if their number and all required constants
are uniform.
\end{lemma}

\begin{proof}
Fix $F$, and take $R_F$ from
Lemma~\ref{probability:thin-geometry}.  We divide the selected
values into nearby collisions, which have a lighter tail, and one
separated set, to which the collection bound applies.

\par\smallskip\noindent\emph{Collision values.}\quad
Define the local field
\[
 D_z=Y_z\mathbf1_{\{\exists w\in\Lambda:\ 0<|w-z|\le R_F,
                                               \ Y_w\ge Y_z\}}.
\]
If $D_z>t$, both $Y_z$ and one distinct neighboring value exceed
$t$.  Independence and Lemma~\ref{probability:weak-two} give
\begin{equation}
 \PP(D_z>t)
 \le\sum_{0<|w-z|\le R_F}\PP(Y_z>t,Y_w>t)
 \le C_F\min\{1,t^{-4}\}.
 \label{probability:collision-four}
\end{equation}
This argument includes laws with atoms.  The dependence range of
$D$ is fixed once $F$ is fixed.  The path budget and
Corollary~\ref{probability:integrated-tails} therefore bound the
expected supremum of its selected tail sums by $C_F(N+1)/T$.

\par\smallskip\noindent\emph{Separated values.}\quad
For any admissible sequence, keep one occurrence of each site in
its original order and pay a factor $m_0$ for multiplicity.  Among
these sites with $Y_z>T$, process values in decreasing order,
breaking ties by a fixed deterministic order.  Retain a site if
its distance from every previously retained site is greater than
$R_F$.  Every discarded site has a distinct retained neighbor
with at least its value, so its value equals $D_z$.  The retained
sites are pairwise separated by more than $R_F$.

Put the retained sites back in their original order.  The all-pairs
order condition still holds, so their images under the deterministic
bijection $\psi_F$ form an admissible directed chain.  They lie
in a deterministic square with radius
\[
 L_F=\max\{1,C(N/F+F\delta N)+C_F\}.
\]
The transported field
$\widetilde Y_w=Y_{\psi_F^{-1}(w)}$ has exactly the original iid
law.  Thus the retained reward is at most the corresponding
$G_n(L_F,x_F)$ in this transported field, for a deterministic
center $x_F$.

Crucially, the same square and the same field permutation work for
every original admissible sequence.  Hence this domination holds
after taking the supremum over those sequences, even though the
pruning and the retained chain depend on the observations.  By
\eqref{probability:collection-assumption}, its expectation is at
most
\[
 C(N/F+F\delta N)+C_F\le CN/F+C_F
 \qquad(\delta\le F^{-2}).
\]
Add the collision estimate and absorb the fixed factor $m_0$ to
obtain \eqref{probability:thin-tail-bound}.

Choose $F$ first so that $C/F<\varepsilon/3$, then choose $T_0$
so that $C_F/T_0<\varepsilon/3$, and set $\delta_0=F^{-2}$.
The remaining terms vanish after division by $N(n)\to\infty$,
uniformly over the laws.  For finitely many types, bound the
supremum of a sum by the sum of the suprema.  For a finite-range
field, first split into finitely many sublattices on which each
fixed template has an iid law, provided the collection hypothesis
holds for every such law.  Apply the argument separately and sum.
\end{proof}

Only one retained chain is charged to the collection bound.  No
packing factor depending on $F$ multiplies $N/F$; all costs of
the increasing collision radius enter $C_F/T$ or the additive
$C_F$.  This is why the order of choices $F$, then $T$, then a
sufficiently large $N$ is effective.

\subsection{Concentration along slowly varying rows}

\begin{lemma}[Uniform row concentration, including omissions]
\label{probability:rows}
For each $n$, let $W^{(n)}_{i,j}$, $i,j\in\Z$, be real random
variables with $|W^{(n)}_{i,j}|\le B$.  Suppose that for every
deterministic sequence $(j_i)$, the variables
$(W^{(n)}_{i,j_i})$ in each residue class $i\pmod q$ are mutually
independent.  The integer $q\ge1$ and $B<\infty$ are fixed.
One sufficient condition is mutual independence of the sigma-fields
generated by entire columns in any set of columns separated by at
least $q$.

Fix $d<\infty$ and $\varepsilon>0$.  Consider rows with
$|j_1|\le N^d$ and
$\sum_{i=1}^{N-1}|j_{i+1}-j_i|\le\delta N$, and omission
sets $O\subset\{1,\ldots,N\}$ with $|O|\le\eta N$.
For sufficiently small fixed $\delta,\eta>0$, depending only on
$\varepsilon,B,q,d$, uniformly in $n$,
\begin{equation}
 \PP\left(
  \sup_{(j_i),O}\left|
    \sum_{i\notin O}\big(W^{(n)}_{i,j_i}
                         -\E W^{(n)}_{i,j_i}\big)
                        \right|>\varepsilon N\right)
 \longrightarrow0\qquad(N\to\infty).
 \label{probability:row-concentration}
\end{equation}
The convergence is exponential up to a polynomial factor.  If
all means equal $m_n$, the comparison sum is $m_n(N-|O|)$.
The conclusion applies to indicators and to bounded sums over any
fixed finite list of offsets or types that separately satisfy the
hypotheses, with adjusted constants.  A direct application to the
sum requires independence at a spacing $q$ appropriate to its
combined input supports.
\end{lemma}

\begin{proof}
The case $B=0$ is immediate.  Otherwise fix a deterministic row
and omission set.  Split the centered sum into its $q$ residue
classes.  Within each class the summands are independent, have
support in intervals of length at most $2B$, and number at most
$N/q+1$.  For the total sum to have absolute value greater than
$\varepsilon N$, some class sum must have absolute value greater
than $\varepsilon N/q$.  Hoeffding's inequality
\cite[Theorem~2, Equation~(2.6)]{Hoeffding1963} and a union bound
therefore give, for $N\ge q$,
\begin{equation}
            2q\exp\left(-\frac{\varepsilon^2N}{4qB^2}\right)
 \label{probability:fixed-row-hoeffding}
\end{equation}
as an upper bound on this probability.

It remains to count the choices over which the supremum is taken.
Put $V=\lfloor\delta N\rfloor$.  For a fixed initial row, stars
and bars gives $\binom{N-1+V}{V}$ choices of the absolute
increments with sum at most $V$.  Their signs have at most $2^V$
choices.  The total number of rows is therefore at most
\[
 (2\lfloor N^d\rfloor+1)2^V\binom{N-1+V}{V}
 \le (2\lfloor N^d\rfloor+1)
                 \exp\{Nh_{\mathrm{row}}(\delta)\},
\]
where
\[
 h_{\mathrm{row}}(\delta)
   =\delta\log2+(1+\delta)\log(1+\delta)-\delta\log\delta
       \longrightarrow0\quad(\delta\downarrow0).
\]
For $0<\eta\le1/2$, there are at most
$(N+1)\exp\{Nh_{\mathrm{bin}}(\eta)\}$ omission sets, where
$h_{\mathrm{bin}}(\eta)=-\eta\log\eta-(1-\eta)\log(1-\eta)$.
Choose $\delta,\eta$ so that the two entropy terms sum to less
than $\varepsilon^2/(8qB^2)$.  Multiplying the counts by
\eqref{probability:fixed-row-hoeffding} gives a polynomial factor
times $2q\exp\{-\varepsilon^2N/(8qB^2)\}$, uniformly in $n$.

If rows are specified only at nonomitted columns, assume the same
bound on their first specified row and on their successive jumps.
Extend to the full interval by holding that first row on the
initial gap and, at each later gap, holding the old row until the
next specified column.  Hold the last row on the final gap.  This
does not increase the total variation and gives a row counted
above.  Indicators have $B=1$.  For a fixed finite sum, apply the
result to each summand with a divided tolerance, then use the
triangle inequality and a finite union bound.  Alternatively, the
sum itself has a fixed bound and can be treated directly when its
combined supports satisfy the column independence condition.
\end{proof}

The supremum in \eqref{probability:row-concentration} is taken
before any row or omission set is selected from the environment.
Its hypotheses must also be checked at the level where it is used.
After conditioning on a shared phase, the same proof is available
if the bound and independence spacing remain uniform, but the
centering is then by conditional means.  Replacing those means by
unconditional ones requires a separate equality or error estimate.
Alternatively, deterministic phase templates or a local phase
construction may directly give the stated unconditional
independence and row invariance.  Averaging functions over finitely
many mesh conventions is permitted when their input supports still
give the required independence spacing.

\subsection{A common dyadic truncation level}

The final observation chooses a truncation from a fixed finite
list.  A first-moment bound is enough for this choice, although it
does not give uniform integrability of the original family.

\begin{lemma}[Good dyadic levels]
\label{probability:dyadic}
Let $U_1^{(n)},\ldots,U_k^{(n)}$ be nonnegative variables with
\[
                 \sup_n\sum_{a=1}^k\E U_a^{(n)}\le A.
\]
Fix $M_0>0$ and an integer $J\ge1$.  For each $n$, there is
$M_n\in\{M_0,2M_0,\ldots,2^{J-1}M_0\}$ such that
\begin{equation}
        M_n\sum_{a=1}^k\PP(U_a^{(n)}>M_n)
                                      \le\frac{2A}{J}.
 \label{probability:good-level}
\end{equation}
Thus this quantity can be made arbitrarily small within a fixed
finite range of levels, uniformly over the laws.  The level may
depend on the law, but is deterministic for that law and is not
chosen from the realization of the field.
\end{lemma}

\begin{proof}
For every $u\ge0$, a geometric sum gives
\[
 \sum_{j=0}^{J-1}2^jM_0\mathbf1_{\{u>2^jM_0\}}\le2u.
\]
Take expectations and sum over $a$.  At least one of the $J$
resulting nonnegative summands is no larger than their average,
which is at most $2A/J$.
\end{proof}

For finitely many thresholds $c_aM$ with fixed $c_a>0$, apply
the lemma to $U_a/c_a$.  The conclusion controls the tail
probability multiplied by its threshold; it makes no assertion
that $\E[U_a^{(n)}\mathbf1_{\{U_a^{(n)}>M\}}]$ is uniformly
small.  When a concentration or comparison estimate is needed at
$M_n$, establish it simultaneously for the finite candidate list,
or use its largest bound $2^{J-1}M_0$ in
Lemma~\ref{probability:rows}.  All concentration tolerances can
then be fixed before $n$ varies.
